\section{Introduction}

Biosafety guidance in the United States places the obligation to assess risk on
the individual laboratory. The sixth edition of \textit{Biosafety in
Microbiological and Biomedical Laboratories} and the fourth edition of the WHO
\textit{Laboratory Biosafety Manual} both frame risk assessment as a local,
procedure-aware activity rather than a matter of agent classification alone
\cite{bmbl6,wholbm4}. ISO 35001:2019 formalises the same expectation as a
management-system requirement, obliging organisations to identify hazards,
assess risk, and mitigate it as a documented cycle \cite{iso35001}.

The obligation is clear. The starting material is not. A laboratory director
who sits down to assess the biorisk of their own bench has no public inventory
to begin from: no list of what the steps of routine testing are, no statement
of what is known about aerosol release from each, and no indication of whether
a given recommendation rests on measurement, on an outbreak investigation, or
on custom. Agent-based risk group classifications describe organisms.
Facility-level containment frameworks describe buildings. Neither describes the
dozen or so physical actions that make up a working day at a diagnostic bench,
and it is those actions that release material.

The scale of the gap is substantial. The Centers for Medicare \& Medicaid
Services record 301,284 laboratories registered under the Clinical Laboratory
Improvement Amendments as of September 2026, of which 15,214 hold a certificate
of compliance and 15,542 a certificate of accreditation \cite{cms2026}. Those
30,756 non-waived settings perform the moderate- and high-complexity testing in
which the steps considered here are routine. Each is required to assess its own
procedures, and each begins from a blank form.

\subsection{Purpose and scope of this brief}

This brief is an inventory and an evidence statement, not a method paper. It
sets out seven bench steps common to non-waived diagnostic testing, defines each
by its boundaries, characterises what in the step is capable of generating an
aerosol, states what has been measured or observed, grades the strength of that
evidence, and records the control set the evidence supports.

The scoring machinery referred to throughout, including the anchored items, the
composition formula and the mapping to ISO 35001:2019 clauses, is specified in
the companion instrument and is not restated here \cite{oyeniyi2026}. Scores
appear in this brief only where they help to order the inventory or to express
the distance between what a step demands and what is known about it.

Three things distinguish an inventory of this kind from existing guidance.

First, it is explicit about evidence. Where guidance exists at procedure level
it rarely distinguishes a recommendation grounded in measurement from one
grounded in precedent. The review by Ravnholdt and colleagues makes the point
directly for aerosol-generating procedures, observing that while some guidance
is informed by empirical data, a significant gap remains in the evidence
supporting it, and that highly automated laboratory equipment in particular has
not been characterised \cite{ravnholdt2025}. Conflating measured and customary
recommendations is not a neutral simplification: it makes a thin evidence base
appear uniform.

Second, it assembles the measurements. The airborne concentration data for
clinical bench work exist in a handful of papers that do not cite one another
and were not designed to be compared. Putting them on one axis, with their
detection limits and their differing conditions visible, is the only way to see
how little there is.

Third, it is built to be used by a survey. A programme surveying laboratory
practice needs to know which questions discriminate between a laboratory that is
safe and one that is merely compliant, and those are the questions that sit on
the evidence gaps. Section~\ref{sec:survey} derives fourteen such items.

\subsection{Relationship to tuberculosis laboratory practice}

Tuberculosis work motivates part of this brief and is treated separately in
Section~\ref{sec:tb}. Three features make it the clearest case in the
inventory. The specimen is respiratory and may carry a high organism burden.
The organism is transmissible at low inhaled dose, so consequence is high
whenever release occurs. And the pathway is long: a single sputum specimen may
pass through reception, aliquoting, molecular testing, digestion and
centrifugation, culture, and identification, with each step adding its own
release opportunity.

Tuberculosis work is also the one place in the inventory where a direct,
quantitative, clinically realistic measurement exists for a step that replaced
another step, which permits the only before-and-after comparison the evidence
base currently supports \cite{banada2010}.
